\section{A certified complex-strip transfer route to the genuine $\Xi$ zeros}
\label{sec:complexstrip}
The finite theta approximation is useful at complex
frequencies as well as on the real axis. In the
Fourier coordinate
\[
s=\frac12+\frac{iz}{2},
\qquad
\Re s=\frac{1-\Im z}{2},
\]
all possible nontrivial off-critical zeros of
$\xi(s)$ lie in the strip $0<|\Im z|<1$.
We now extend the exact square-lattice tail
majorant far enough into the complex plane
to cover \emph{the entire zero-relevant strip}.

\begin{theorem}[Global horizontal-strip
theta approximation]\label{thm:complexstrip}
Let $F(z)=\xi(1/2+iz/2)/4$ and let $F_M(z)$
be the incomplete-gamma finite-mode transform
defined in Section~\ref{sec:squarelattice}.
Put $q=\pi(M+1)^2$ and
$E_M=3qe^{-q}$. Then for every integer
$k\ge0$,
\begin{equation}
\boxed{
\sup_{\substack{z\in\mathbb C\\|\Im z|\le3}}
|F^{(k)}(z)-F_M^{(k)}(z)|
\le\frac{k!}{4^k}E_M.
}
\label{eq:stripderivative}
\end{equation}
For $k=0$ the larger strip
$|\Im z|\le7$ satisfies
\begin{equation}
\boxed{
\sup_{|\Im z|\le7}|F(z)-F_M(z)|
\le E_M.
}
\label{eq:stripzero}
\end{equation}
For each fixed $n$ the real-frequency
Laguerre difference bound from
Theorem~\ref{thm:uniformFourier}
also holds for complex $z$ with
$|\Im z|\le3$, interpreting the
Laguerre expression as its analytic
polynomial in $F(z)$ and its derivatives.
\end{theorem}
\begin{proof}
For $z=x+iy$ and real $r\ge0$, both
$|\cos(zr)|$ and $|\sin(zr)|$ are
at most $e^{|y|r}$. As before,
$r^k\le k!e^{4r}/4^k$.
Using $\phi_m(r)\le
2q_m^2e^{9r-q_me^{4r}}$,
the omitted $k$th derivative obeys
\begin{align}
|F^{(k)}(z)-F_M^{(k)}(z)|
&\le \frac{4k!}{4^k}
\sum_{m>M}q_m^2
\int_0^\infty
e^{(13+|y|)r-q_me^{4r}}\,dr
\nonumber\\
&=\frac{k!}{4^k}
\sum_{m>M}q_m^2
\int_1^\infty
t^{(9+|y|)/4}e^{-q_mt}\,dt.
\label{eq:complexmajorant}
\end{align}
If $|y|\le3$, the exponent of $t$
is at most $3$.
The same incomplete-gamma
integer-degree calculation as in
\eqref{eq:tailterm} therefore bounds
the sum by $E_M$.
For $k=0$ we omit the auxiliary
$e^{4r}$, giving instead the exponent
$(5+|y|)/4\le3$ whenever $|y|\le7$.
This proves the two stated majorants.

The absolute derivative bound
$|F^{(j)}(z)|,|F_M^{(j)}(z)|
\le Cj!/4^j$ follows by applying
the same argument to the full sum.
The product comparison
and binomial-factorial cancellation
used to prove \eqref{eq:uniformlag}
work identically with complex
absolute values on $|\Im z|\le3$,
proving the last assertion.
\end{proof}

\begin{corollary}[Effective Rouch\'e zero-count transfer]
\label{cor:rouche}
Let $\Omega$ be a bounded Jordan
domain contained in $|\Im z|<7$
whose closure lies in that strip.
If
\begin{equation}
\boxed{
\inf_{z\in\partial\Omega}|F_M(z)|>E_M,
}
\label{eq:rouchegap}
\end{equation}
then $F$ and $F_M$ have precisely
the same number of zeros in
$\Omega$, counted with multiplicity.
\end{corollary}
\begin{proof}
On $\partial\Omega$,
Theorem~\ref{thm:complexstrip}
implies $|F-F_M|\le E_M<|F_M|$.
The classical Rouch\'e theorem
therefore applies.
\end{proof}

\paragraph{Numerical strip audit.}
For $M=2$, the uniform bound is
$E_2=3(9\pi)e^{-9\pi}
\simeq4.4578603\cdot10^{-11}$.
An independent 55-digit
incomplete-gamma and
completed-zeta replay gave:
\begin{center}\small
\begin{tabular}{c c c}
\hline
$z$& $|F_2(z)|$ & $|F(z)-F_2(z)|$\\
\hline
$28.269+0.8i$&$1.3904\,10^{-4}$&
$1.3869\,10^{-11}$\\
$39.5+0.99i$&$1.4461\,10^{-5}$&
$1.3131\,10^{-11}$\\
$100+0.5i$&$8.1041\,10^{-12}$&
$8.1050\,10^{-12}$\\
\hline
\end{tabular}
\end{center}
These are numerical point samples,
not certified minimum moduli
over closed contours. In particular,
although \eqref{eq:stripzero}
holds uniformly in $\Re z$,
the $M=2$ approximant at
$100+0.5i$ is already smaller
than its rigorous absolute
error allowance; \eqref{eq:rouchegap}
cannot be inferred there.

\paragraph{A concrete first-zero contour candidate.}
In a further numerical (not interval-certified)
experiment, take the circle
$|z-28.269|=1/2$ and $M=2$.
At $96$ equally spaced points the
minimum sampled modulus is approximately
$7.45452\cdot10^{-5}$, about
$1.67\cdot10^6$ times the proven
strip-tail bound $E_2$.
The unwrapped sampled argument of
$F_2$ winds once around the origin.
Radii $1/4$ and $1$ give the same
numerical winding number and
sample margins above $8.9\cdot10^5E_2$.
This is a promising target for
a finite certified zero count,
but \emph{a discrete sample minimum is
not a lower bound on the entire
continuous contour}. The exact
missing certification is a
directed-rounding enclosure of all
sample values joined to an analytic
modulus-of-continuity bound (or
a uniform analytic contour estimate).
No rigorous winding/zero count
is claimed from these samples.
The experiment is reproduced by
\texttt{certificates/xi\_first\_zero\_rouche\_contour\_numerical.py}.

\paragraph{A finite 4,096-node certification design.}
The above exploratory contour admits a
particularly simple rigorous
modulus-of-continuity criterion.
The exact derivative bound of
Theorem~\ref{thm:complexstrip}
gives, throughout $|\Im z|\le3$,
\[
|F_2'(z)|\le C/4<1/8.
\]
Indeed $3<\pi<22/7$,
$e^3>20$, and the explicit
value of $C$ imply the
purely rational check
\[
C<
\frac{22/7+3+2+2/3}{20}
+\frac{12(22/7)}{20^4}<\frac12.
\]
With $N=4096$ equally spaced
points on the radius-$1/2$
circle, every point of the
circle lies at distance at most
$\pi/(2N)$ from a sample.
Hence the Lipschitz loss
from the nearest sample is
bounded by
\[
\frac18\frac{\pi}{2N}
<\frac{22}{7\cdot65536}.
\]
Also the rigorous $M=2$ strip
error satisfies
\[
E_2<\frac{27(22/7)}{20^9}
<10^{-9}.
\]
Therefore the following is an
\emph{exact finite certificate target}:
if all 4,096 sampled values of
$F_2$ can be enclosed by directed
complex interval arithmetic so
that
\[
|F_2(z_j)|\ge7\cdot10^{-5}
\qquad(0\le j<4096),
\]
then necessarily on the
\emph{entire continuous contour}
\[
|F_2(z)|>
7\cdot10^{-5}
-\frac{22}{7\cdot65536}
>E_2.
\]
The inequality and exact rational
thresholds are kernel-checked in
\texttt{XiKernel/ContourMeshGate.lean}.
A validated winding number of
the 4,096-node polygon would then
produce a rigorous single-region
Rouch\'e certificate.
The necessary 4,096 directed
complex enclosures and winding
certificate have \emph{not}
yet been constructed: existing
96-node mpmath samples are
numerical only. This is a
finite proof obligation,
not a completed zero-count
certificate.

\subsubsection{An interval-certified first-zero disk}
\label{subsec:firstzerocert}
The preceding 4,096-node design has a simpler
replacement: one can compare $F_2$ with its
\emph{linear Taylor polynomial} on the circle.
This requires just three Taylor coefficients
and one strictly positive remainder integral,
rather than thousands of validated contour
values or an independent winding computation.

\begin{theorem}[A finite certified zero count for
the genuine $\Xi$ transform]\label{thm:firstzerodisk}
For $z_0=28269/1000$, the entire function
\[
F(z)=\frac14\xi\!\left(\frac12+\frac{iz}{2}\right)
\]
has exactly one zero, counted with
multiplicity, in the disk
\begin{equation}
\boxed{|z-z_0|<\frac12.}
\label{eq:firstzerodisk}
\end{equation}
That unique zero is real and simple.
The proof is a finite interval-arithmetic
certificate with explicit Simpson and
theta-tail error bounds, not a proof of RH.
\end{theorem}

\begin{proof}
Take the two-mode positive-axis Fourier
transform
\[
F_2(z)=2\int_0^\infty
\bigl(\phi_1(r)+\phi_2(r)\bigr)\cos(zr)\,dr
\]
and write $w=z-z_0$. Taylor's formula gives
\[
F_2(z_0+w)=a_0+a_1w+a_2w^2+R_3(w),
\]
where $a_k=F_2^{(k)}(z_0)/k!$.
For $|w|\le1/2$, the integral remainder
obeys
\begin{equation}
|R_3(w)|\le
\frac1{24}\int_0^\infty
\bigl(\phi_1(r)+\phi_2(r)\bigr)
r^3e^{r/2}\,dr .
\label{eq:firstzeroR}
\end{equation}
The four real quantities on the right
and in the Taylor coefficients are
enclosed by outward-rounded interval
Simpson quadrature on $[0,2]$,
with analytic bounds for the omitted
$r>2$ part.

For completeness, the verified
uniform mode-derivative constants
for $0\le j\le4$ are
\[
(B_0,B_1,B_2,B_3,B_4)=
(30,1272,79320,6800400,762395760).
\]
Indeed each mode derivative has the form
$q^{-1/4}v^{5/4}e^{-v}P_j(v)$,
with $v=qe^{4r}\ge3$ and
$|P_j(v)|\le\|P_j\|_1v^{j+1}$.
The elementary exponential comparison
$v^{j+3}e^{-v}\le(j+3)!$
yields $B_j=\|P_j\|_1(j+3)!$.
Applying Leibniz's rule to
$2(\phi_1+\phi_2)r^k\cos(z_0r)/k!$
(or the corresponding odd derivative)
provides valid fourth-derivative bounds
for $k=0,1,2$, respectively,
\[
8.388\cdot10^9,\quad
1.706\cdot10^{10},\quad
1.735\cdot10^{10}.
\]
For the positive remainder integrand
in \eqref{eq:firstzeroR}, a valid
fourth-derivative bound is
$1.637\cdot10^9$.
All constants are rounded
\emph{upward} here; the code uses
their exact rational values.
With $8192$ equal Simpson panels
on $[0,2]$, the formula
\[
|\mathrm{error}|
\le\frac2{180}
\left(\frac2{8192}\right)^4
\sup_{[0,2]}|g^{(4)}|
\]
bounds the four quadrature errors
by $3.32\cdot10^{-7}$,
$6.74\cdot10^{-7}$,
$6.85\cdot10^{-7}$ and
$6.47\cdot10^{-8}$.

The remaining $r>2$ tails
are less than $10^{-20}$.
For this last inequality, use
$1000<e^8<5000$, $3<q_m<16$
for $m=1,2$, $e^3>20$, and
the elementary integer-degree
incomplete-gamma upper bound
\[
e^{-q_me^8}
\left(q_me^{24}+3e^{16}
+\frac{6e^8}{q_m}+\frac6{q_m^2}\right)
<10^{14}20^{-1000}.
\]
These rational comparisons are
checked explicitly by positive
Taylor sums, including an upper
tail estimate for $e^2<8$.

The executable certificate therefore
establishes the \emph{strict rational}
enclosures
\begin{equation}
\boxed{
\begin{gathered}
|a_0|<10^{-6},\quad
a_1<-17\cdot10^{-5},\\
|a_2|<52\cdot10^{-6},\quad
|R_3(w)|<6\cdot10^{-6}
\quad(|w|\le\tfrac12).
\end{gathered}}
\label{eq:firstzeroTaylor}
\end{equation}
As an independent numerical-resolution
test, the certificate was repeated
using $16384$ Simpson panels; both
full directed-rounding runs establish
the same rational inequalities.

On the circle $|w|=1/2$,
\eqref{eq:firstzeroTaylor} gives
\begin{align}
|F_2(z)|
&\ge|a_1w|-|a_0|-|a_2w^2|-|R_3(w)|\nonumber\\
&> \frac{17\cdot10^{-5}}2
-10^{-6}-\frac{52\cdot10^{-6}}4
-6\cdot10^{-6}\nonumber\\
&=65\cdot10^{-6}.
\label{eq:linearmargin}
\end{align}
Therefore $F_2$ has exactly one
zero in the disk by Rouch\'e's
theorem, compared with
$a_1(z-z_0)$.

The genuine theta-tail inequality
\eqref{eq:stripzero} yields, throughout
the same disk and its boundary,
\[
|F(z)-F_2(z)|
\le27\pi e^{-9\pi}
<\frac{27(22/7)}{20^9}
<10^{-9}.
\]
A second application of Rouch\'e
shows that $F$ has precisely one
zero there. Since $F(\bar z)=
\overline{F(z)}$ and the disk is
invariant under conjugation,
a nonreal zero would occur with
its distinct conjugate, impossible
for a single zero counted with
multiplicity. Hence the unique
zero is real and simple.
\end{proof}

\paragraph{Trust boundary and significance.}
The stand-alone checker
\texttt{certificates/xi\_first\_zero\_interval.py}
uses \texttt{mpmath.iv} version 1.3.0
to enclose $\pi$, exponentials, sine
and cosine at the Simpson nodes;
every interval endpoint is converted
to an exact rational number before
the final inequalities are decided.
The derivative majorants,
composite Simpson remainder,
integral tails and final Rouch\'e
margin are explicit rational
inequalities. The implementation
was replayed at both $8192$ and
$16384$ panels. The proof relies
on the correctness of the interval
library's transcendental enclosures
and standard analytic theorems.
The corresponding
\texttt{XiKernel/FirstZeroDiskGate.lean}
checks the final rational gap but
does not formalize the transcendental
quadrature, Fourier identity or
Rouch\'e theorem end-to-end.
A single certified classical
real zero is a useful reproducibility
benchmark; it does not verify
infinitely many critical zeros
or establish RH, and no novelty
is asserted for this zero's existence.

\begin{proposition}[The exact remaining
global zero-migration criterion]
\label{prop:zeromigration}
For each $0<\epsilon<1/2$ and
$R>0$, define
\[
K_{\epsilon,R}=
\{z:|\Re z|\le R,\quad
\epsilon\le|\Im z|\le1-\epsilon\}.
\]
The Riemann hypothesis is equivalent
to the following \emph{eventual
finite-cutoff zero evacuation} property:
for every such pair $(\epsilon,R)$,
there is an $M_0$ for which
\[
F_M(z)\ne0\quad
\text{for all }z\in K_{\epsilon,R}
\text{ and every }M\ge M_0.
\]
\end{proposition}
\begin{proof}
If RH holds, $F$ has no zeros
on the compact set $K_{\epsilon,R}$.
Its minimum modulus there is
strictly positive. The uniform
error \eqref{eq:stripzero}
tends to zero as $M\to\infty$,
so all sufficiently large $F_M$
are zero-free on this compact set.

Conversely, suppose $F$ has an
off-real zero $z_0$ in
$|\Im z_0|<1$. Choose
$\epsilon,R$ such that $z_0$
lies in the interior of
$K_{\epsilon,R}$. Uniform
holomorphic convergence
$F_M\to F$ on a disk about
$z_0$, together with the fact
that $F$ is not identically
zero, implies by Hurwitz's
theorem (or the argument
principle on a small zero-free
contour) that every
sufficiently large $F_M$
has a zero in that disk.
This contradicts eventual
zero evacuation. The
classical correspondence
between off-critical
zeta zeros and off-real
$\Xi$ zeros completes the
equivalence.
\end{proof}

\paragraph{The substantive unresolved step.}
Proposition~\ref{prop:zeromigration}
is a mathematically correct
\emph{reformulation}, not a solution.
The new effective estimate handles
each compact region but supplies
only an \emph{absolute} error.
For each fixed $M$, $F_M(z)$
tends to zero as
$|\Re z|\to\infty$ in
bounded horizontal strips;
therefore the required
positive contour-modulus gap
cannot hold uniformly at
arbitrarily high frequencies
with a fixed $E_M>0$.
Any actual RH proof by this
route must provide an
adaptive cutoff $M=M(R)$,
certified contour moduli
or another global-in-frequency
zero-exclusion estimate.
A finite number of successful
Rouch\'e checks does not
establish the full statement.

\paragraph{Reproducibility and status.}
The exact rational sublemmas
on complex-strip exponents
and convolution coefficients
are verified by
\texttt{certificates/xi\_complex\_strip\_exact.py}
and
\texttt{XiKernel/ComplexStripTransfer.lean}.
An independent numerical
replay is in
\texttt{certificates/xi\_complex\_strip\_numerical.py}.
The Rouch\'e/Hurwitz transfer
is proved analytically here;
it is not yet an end-to-end
Lean proof of the genuine
Jacobi-theta or $\Xi$ functions.
The criterion is classical
complex analysis combined with
the explicit arithmetic
theta majorant; no priority
claim is made, and RH remains
unresolved.
